\begin{proof}

\begin{enumerate}[label={$\langle1\rangle$\arabic*.},leftmargin=*]
\item We extend the boundary argument of P4.20 in the
precise form used here. Suppose a deformation retains
the smooth critical curve near a fixed transverse
boundary sphere and retains transverse type $A_2$.
Near a boundary point, the holomorphic Morse lemma
in the nondegenerate transverse variable followed
by extraction of a cube root of a unit gives
coordinates with defining function $a^2+b^3$.
The third coordinate parametrizes the critical
curve. Its transversality to the boundary lets
one real third-coordinate variable be solved
from the boundary equation; the complex
coordinates $a,b$ remain free. Thus the
restriction to the boundary is submersive
at every nonzero level. Compactness and
the smooth part of the boundary give the
same uniform nonzero-level submersion as
in P4.20. Its finite-critical-value and
proper Ehresmann transport proof applies
unchanged. A compact spherical class of
nonzero square in a deformed nearby
fibre would therefore lie in an actual
nearby torus fibre, which is impossible.
All deformations below retain this
boundary transverse type.

\item Choose holomorphic coordinates $(x,y,z)$
so the smooth critical curve is $x=y=0$.
The transverse Hessian along it is a
symmetric rank-one matrix over the
discrete valuation ring $\mathbb C\{z\}$.
It is not the zero matrix generically.
It can be diagonalized holomorphically
to $h(z)x^2$. Here is the elementary
ring calculation. The determinant is
zero. Choose a diagonal entry of
minimal order; one exists since the
square of the off-diagonal entry is
the product of the two diagonal entries.
The ratio of the off-diagonal entry
to the chosen entry is holomorphic.
Completing its square gives the desired
invertible coordinate change. The
coefficient $h$ is not identically zero.

\item First suppose that the Hessian has
rank one at $p$. The Morse lemma with
parameters $(y,z)$ writes the reduced
equation as
\[
 g=X^2+y^3h_1(y,z).
\]
Indeed after removing the $X$ direction,
the remaining function and its first
two $y$ derivatives vanish on the
critical curve; the third derivative
is not identically zero there. If
$h_1(0,0)\ne0$, extracting its cube
root in the $y$ coordinate proves the
asserted product form.

\item Otherwise $h_1(0,z)$ has a zero of
finite positive order. Use
$g_\eta=X^2+y^3(h_1+\eta)$.
For arbitrarily small generic $\eta$,
the function $h_1(0,z)+\eta$ has a
simple zero $z_\eta$ near zero.
This follows directly from one-variable
Weierstrass preparation, avoiding the
finite critical values of $h_1(0,z)$.
At $(0,0,z_\eta)$ the function
$Z=h_1(y,z)+\eta$ is a coordinate,
since its $z$ derivative is nonzero.
The actual germ is
$X^2+y^3Z$, which is excluded by
P4.16: its nearby fibre has a
spherical class of square $-6$.
On a fixed sufficiently small boundary,
$h_1(0,z)$ is nonzero at the critical
curve, and this persists for small
$\eta$. The deformed transverse type
is therefore still $A_2$ there.
The preceding boundary transport gives
the contradiction.

\item It remains to exclude Hessian rank
zero at $p$, so $h(0)=0$. In the
diagonalizing coordinates add
$\varepsilon y^3$ to $g$. The
critical curve and its transverse
Hessian do not change. The cubic
coefficient in the transverse kernel
remains nonzero near the original
compact boundary for small
$\varepsilon$, so this deformation
has the boundary property just proved.
Choose $\varepsilon$ so that the
coefficient of $y^3$ at the origin
is nonzero.

\item On a smaller neighbourhood, Weierstrass
preparation in $y$ gives a unit times
\[
 P=y^3+a(x,z)y^2+b(x,z)y+c(x,z).
\]
Restriction to $x=0$ is a unit times
$y^3$, so all three coefficients
vanish at $x=0$. The critical-curve
equation $g_x(0,0,z)=0$ first gives
$c_x(0,z)=0$. In the identity
$g=U P$ the mixed derivative on
the critical curve is then
$g_{xy}=U b_x$, since the terms
containing $c,c_x,b$ vanish there.
The zero $xy$ Hessian entry therefore
gives $b_x(0,z)=0$. Thus
$x^2\mid b$, $x^2\mid c$, and
$x\mid a$.
The Tschirnhaus change $y=Y-a/3$
therefore gives exactly
\[
 P=Y^3+x^2H(x,Y,z).
\]
The coefficient $H(0,0,z)$ is a unit
multiple of the original Hessian
coefficient $h(z)$. It is not
identically zero but vanishes at
zero. Perturb it to $H+\eta$.
At a simple root of
$H(0,0,z)+\eta$, the coordinates
$(x,Y,Z=H+\eta)$ give
$Y^3+x^2Z$. P4.18 excludes this
binomial germ by its compact
spherical class of nonzero square.

\item There are two nested transport
neighbourhoods in this argument.
First fix the original ball for
$g\to g+\varepsilon y^3$.
Then choose a smaller ball in the
Weierstrass coordinate neighbourhood,
where $g+\varepsilon y^3$ has only
the original critical curve and its
boundary has transverse type $A_2$.
The second perturbation transports
its spherical class into a nearby
fibre on this smaller ball. Inclusion
puts it in a fibre on the original
ball; the first boundary transport
puts it in a fibre of $g$.
This does not require extending the
second perturbation outside its
coordinate neighbourhood.
Units in the reduced equations are
retained during the perturbations.
At a binomial special point a unit
can be absorbed by the positive
weighted coordinate scaling of
P4.16--P4.18.

\item We have ruled out both a zero
Hessian and a vanishing cubic
coefficient. The product form
follows. Its normalization is
the finite map
\[
 (u,z)\longmapsto (a=u^3,b=-u^2,z),
\]
whose source is smooth, whose
inverse rational parameter off
$b=0$ is $u=-a/b$, and whose
restriction to the conductor is
the identity on the $z$ coordinate.

\end{enumerate}
\end{proof}
